\documentclass[11pt,a4paper,fleqn]{article}
\usepackage[T1]{fontenc}
\usepackage[utf8]{inputenc}
\usepackage[english,ngerman]{babel}
\usepackage{lmodern}
\usepackage[margin=2.3cm]{geometry}
\usepackage{amsmath,amssymb,bm}
\usepackage{xcolor}
\usepackage{tikz}
\usetikzlibrary{arrows.meta,calc,positioning,shapes.geometric,babel}
\usepackage{pgfplots}
\usepackage{pgfplotstable}
\pgfplotsset{compat=1.18}
\usepackage{tcolorbox}
\usepackage{booktabs}
\usepackage{array}
\usepackage{tabularx}
\usepackage{multicol}
\usepackage{enumitem}
\usepackage{etoolbox}
\usepackage[font=small,labelfont=bf,hypcap=false]{caption}
\usepackage[colorlinks=true,linkcolor=black,citecolor=blue!50!black,urlcolor=blue!50!black]{hyperref}
\usepackage{parskip}
\usepackage[compact]{titlesec}
\titlespacing*{\section}{0pt}{2.2ex plus .6ex minus .4ex}{1ex plus .2ex}
\titlespacing*{\subsection}{0pt}{1.6ex plus .5ex minus .3ex}{0.6ex plus .2ex}
\usepackage{fancyhdr}
\setlength{\headheight}{14pt}
\pagestyle{fancy}
\fancyhf{}
\lhead{\small DeCAF --- Bohrer-Tracking: Problem und Lösung (verständlich)}
\rhead{\small Stand 13.11.2026}
\cfoot{\thepage}
\setlist[itemize]{leftmargin=*, itemsep=2pt, parsep=0pt, topsep=2pt}
\setlist[enumerate]{leftmargin=*, itemsep=2pt, parsep=0pt, topsep=2pt}
\emergencystretch=2em
% keine einzelnen Zeilen am Seitenanfang oder -ende
\widowpenalty=10000
\clubpenalty=10000
\hyphenation{kur-zen Roh-auf-nah-men}
% Literaturverzeichnis kompakt und zweispaltig (Ueberschrift ausserhalb der Spalten)
\patchcmd{\thebibliography}{\section*{\refname}}{}{}{}
\apptocmd{\thebibliography}{\setlength{\itemsep}{0.5pt}\setlength{\parsep}{0pt}}{}{}

% Formelbox wie in der Arbeit (Header.tex, mybox)
\newtcolorbox{mybox}[1]{colback=red!5!white,colframe=red!75!black,fonttitle=\bfseries,title={#1}}

% Farben wie im Langdokument; Beschriftungen schwarz
\definecolor{ColA}{RGB}{224,0,0}          % heute (rot)
\colorlet{ColB}{blue!65!black}            % Halten (hier nicht gezeichnet)
\definecolor{ColN}{RGB}{31,163,31}        % neuer Schaetzer (gruen wie die gezeichnete Spitze)
\colorlet{CarryC}{orange!85!black}        % carry-Bilder

% Notation wie in der Arbeit (Methods 3.3.7, 3.4)
\newcommand{\Tcm}{\mathbf{T}_{\mathrm{cm}}}
\newcommand{\Ptd}{\mathbf{P}^{\mathrm{tip}}_{\mathrm{d}}}
\newcommand{\Ptc}{\mathbf{P}^{\mathrm{tip}}_{\mathrm{c}}}
\newcommand{\stip}{\sigma_{\mathrm{tip}}}
\newcommand{\xb}{\bar{\mathbf{x}}_k}
\newcommand{\xh}{\hat{\mathbf{x}}}
\newcommand{\Sb}{\bar{\boldsymbol{\Sigma}}_k}

\begin{document}

\begin{center}
  {\LARGE Bohrer-Tracking: Problem und Lösungsansatz}\\[2mm]
  {\large Warum die gezeichnete Spitze springt und was ein Schätzer mit Gedächtnis daran ändert}\\[2mm]
  {\small DeCAF --- Masterarbeit \quad$\cdot$\quad Stand 13.11.2026}
\end{center}

\begin{mybox}{Kurzfassung}
Heute bestimmt das Tracking die Lage des Bohrers in jedem Kamerabild neu, ohne sie über die Bilder davor zu
mitteln. Je nachdem, ob die mit IR-Punkten verbesserte, \emph{fusionierte} Pose eines Bildes vier Prüfungen
besteht, wird entweder sie gezeichnet oder die reine Farbpose, meist plus einem gemerkten Versatz. Die Farbpose ist
vor allem in der Tiefe unsicher und liegt von Bild zu Bild verschieden weit daneben; der gemerkte Versatz gleicht
das nur für das Bild aus, auf dem er gemerkt wurde. Deshalb springt die Spitze auch bei stillstehendem Bohrer, im
Labortest am 11.11.\ 83-mal um mehr als 5\,mm. Der neue Ansatz ersetzt das Umschalten durch einen einzigen
Schätzer. Er merkt sich die letzte Lage, verwendet in jedem Bild Farbecken und IR-Punkte nach ihrer
Zuverlässigkeit und wertet auffällige Messungen ab, statt die IR-Punkte ganzer Bilder zu verwerfen (außer bei zu
wenigen oder grob unpassenden IR-Punkten). Mit eingespanntem Bohrer zittert die Spitze auf den Testdaten nur noch
etwa ein Drittel so stark. Der Labortest selbst lässt sich ohne Rohbilder nicht nachrechnen. In einer ähnlichen
Aufnahme aus den Abstimmdaten, handgeführt mit dem Etikett am Bildrand, sinken die Sprünge über 5\,mm von 16 auf 0;
die Streuung bleibt dort aber gleich, und im Median liegt die Spitze sogar weiter von der Referenz entfernt
(Abschnitt~\ref{sec:grenzen}). Den bleibenden Abstand der Spitze zur Referenz von etwa 1,5 bis 2\,mm, den es schon
heute gibt, beseitigt er nicht. Der Schätzer ist in die C++-Engine eingebaut und in der Live-Anzeige voreingestellt
an; ausgeschaltet sind Rechnung und Anzeige wie vorher.
\end{mybox}

Herleitungen, alle Einstellungen und die genauen Bedingungen jeder Zahl stehen im ausführlichen Bericht \glqq
Bohrer-Tracking: Problem und neuer Schätzer\grqq{} im selben Ordner.

%==========================================================================================
\section{Was heute passiert}
\label{sec:heute}

Die Kamera liefert in jedem Takt ein Farbbild und, wenn ein Werkzeug geladen ist, ein Bildpaar der beiden
IR-Kameras. Daraus bestimmt die C++-Tracking-Engine die \emph{Pose} des Bohrers, also wo das Zylinder-Etikett im
Raum liegt und wie es gedreht ist. Die Spitze $\mathbf{P}^{\mathrm{tip}}$ ist ein Punkt, der starr am Etikett
hängt, 179,36\,mm von dessen Ursprung entfernt. Pro Bild läuft Folgendes ab (Abb.~\ref{fig:heute}):

\begin{enumerate}
  \item \textbf{Farbpose.} Die Engine findet die Ecken des Schachbrettmusters auf dem Etikett im Farbbild, ordnet
        sie den Modellecken zu und berechnet daraus die \emph{Farbpose}. Ein kleiner Filter je Ecke glättet die
        Eckpositionen leicht; die Pose selbst wird nicht geglättet.
  \item \textbf{IR-Punkte und fusionierte Pose.} Dieselben Ecken werden in beiden IR-Bildern gesucht. Eine Ecke,
        die in beiden gefunden wird, ergibt durch Triangulation (wie beim Sehen mit zwei Augen) einen 3D-Punkt,
        hier \emph{IR-Punkt} genannt. Aus Farbecken und IR-Punkten zusammen berechnet die Engine die
        \emph{fusionierte Pose}. Die Farbpose ist in der Tiefe (entlang der Sichtlinie der Kamera) und in der
        Kippung des Zylinders unsicher. Die IR-Punkte sind dagegen schon räumliche Punkte mit eigener Tiefe; mit
        ihnen ist die fusionierte Pose in diesen Richtungen viel besser bestimmt.
  \item \textbf{Vier Prüfungen.} Vier Ja/Nein-Prüfungen entscheiden für das ganze Bild, ob die fusionierte Pose
        verwendet wird:
        \begin{itemize}[topsep=0pt, partopsep=0pt, itemsep=0pt, parsep=0pt]
          \item mindestens 6 IR-Punkte;
          \item die IR-Punkte passen gut zum Modell (Restfehler höchstens 2,5\,mm);
          \item die fusionierte Pose legt den Ursprung des Etiketts höchstens 6\,mm neben den der Farbpose;
          \item die fusionierte Pose passt im Farbbild nicht wesentlich schlechter als die Farbpose.
        \end{itemize}
  \item \textbf{Zeichnen.} Bei bestandenen Prüfungen wird die fusionierte Pose gezeichnet; solche Bilder heißen
        hier \emph{fused}-Bilder. Die Anzeige merkt sich dabei den Versatz zwischen fusionierter Spitze und
        Farbspitze. Sonst wird die Farbpose gezeichnet, meist plus dem gemerkten Versatz, solange er höchstens
        1\,s alt ist (\emph{carry}-Bilder, \glqq mitgetragen\grqq). Der Versatz ist eingefroren: Er dreht sich
        nicht mit, wenn sich der Bohrer bewegt, und endet nach 1\,s abrupt. Ohne Versatz, etwa direkt nach dem
        \emph{Einrasten} (dem ersten Finden des Etiketts nach dem Start oder nach einem Verlust), bleibt es bei
        der reinen Farbpose.
\end{enumerate}

Die Engine findet die Farbpose, indem sie eine Anfangspose, den \emph{Startwert}, schrittweise verbessert. Der
Startwert kommt immer aus der Farbpose des Bildes davor, nie aus der fusionierten Pose (in Abb.~\ref{fig:heute}
die gestrichelte Schleife). Die fusionierte Pose wirkt also nur auf das, was gezeichnet wird; ihre bessere Tiefe
hilft dem nächsten Bild nicht.

\begin{figure}[!htb]
\centering
\begin{tikzpicture}[
  font=\footnotesize,
  box/.style={draw, rounded corners=1.5pt, align=left, inner sep=2.5pt, fill=white, text width=#1,
              execute at begin node={\hyphenpenalty=10000\exhyphenpenalty=10000}},
  box/.default=26mm,
  dec/.style={draw, diamond, aspect=1.7, align=center, inner sep=0pt, text width=21mm,
              execute at begin node={\hyphenpenalty=10000\exhyphenpenalty=10000}},
  arr/.style={-{Stealth[length=1.8mm]}, semithick},
  fb/.style={-{Stealth[length=1.8mm]}, dashed, black!60},
  lab/.style={font=\footnotesize, fill=white, inner sep=1pt, align=center}
]
\node[box=24mm] (F) at (0,0) {\textbf{1 Farbpose}\\ Ecken im Farbbild finden, daraus die Pose};
\node[box=30mm, right=4mm of F] (I) {\textbf{2 IR-Fusion}\\ IR-Punkte aus beiden IR-Bildern; Pose aus Farbecken und
IR-Punkten};
\node[dec, right=4mm of I] (D) {\textbf{3 Prüfungen}\\ bestanden?};
\node[box=31mm, anchor=west] (FU) at ($(D.east)+(6mm,14mm)$) {\textbf{4 fused}: fusionierte Pose zeichnen, Versatz
zur Farbpose merken};
\node[box=31mm, anchor=west] (CA) at ($(D.east)+(6mm,-14mm)$) {\textbf{4 carry} (meist): Farbpose plus dem gemerkten
Versatz zeichnen};
\node[draw, rounded corners=1.5pt, fill=white, align=center, text width=11mm, inner sep=2.5pt,
      minimum height=38mm] (Z) at ($(FU.east)!0.5!(CA.east)+(11mm,0)$) {Spitze\\ ins Bild};
\draw[arr] (F) -- (I);
\draw[arr] (I) -- (D);
\draw[arr] (D.north) |- node[lab, pos=0.25, left=0.5mm] {ja} (FU.west);
\draw[arr] (D.south) |- node[lab, pos=0.25, left=0.5mm] {nein} (CA.west);
\draw[arr] (FU.east) -- (FU.east -| Z.west);
\draw[arr] (CA.east) -- (CA.east -| Z.west);
\draw[fb] ([xshift=-5mm]FU.south east) -- node[lab, left=0.8mm, align=right, fill=none, text=black] {Versatz\\ eingefroren,\\ höchstens 1\,s} ([xshift=-5mm]CA.north east);
\draw[fb] ([xshift=-7mm]F.south) -- ++(0,-5mm) -| ([xshift=7mm]F.south);
\node[lab, anchor=north, align=center] at ($(F.south)+(0,-5.8mm)$) {Startwert für\\ das nächste Bild};
\end{tikzpicture}
\caption{Heutige Verarbeitung pro Kamerabild, vereinfacht. Je nach Ergebnis der Prüfungen wird entweder die
fusionierte Pose aus Farbecken und IR-Punkten gezeichnet (fused) oder die Farbpose, meist plus dem eingefrorenen
Versatz des letzten fused-Bildes (carry). Ohne gemerkten Versatz wird die reine Farbpose gezeichnet. Der Startwert
für das nächste Bild kommt immer aus der Farbpose; über die Zeit gemittelt wird die Pose nicht.}
\label{fig:heute}
\end{figure}

\textbf{Ein Vergleich aus dem Alltag.} Man denke an zwei Küchenwaagen: Die eine wiegt recht genau, die andere
liegt je nach Lage des Pakets mal mehr, mal weniger daneben. Wer zwischen beiden wechselt und bei der zweiten einen
einmal gemerkten Unterschied aufschlägt, sieht die Anzeige trotzdem springen, denn der gemerkte Unterschied stimmte
nur für das eine Mal. So verhält sich heute die gezeichnete Spitze: Die fusionierte Pose ist die genauere Waage,
die Farbpose die ungenaue, und der carry-Versatz ist der einmal gemerkte Unterschied.

%==========================================================================================
\section{Wo das Problem liegt}
\label{sec:problem}

\textbf{Datengrundlage.} Ausgewertet wurden Aufnahmen mit Rohbildern, die sich offline erneut durch das Tracking
schicken lassen. Meist steckt die Spitze dabei in einem Sitz des Kalibrierbretts. Wo dieser Sitz liegt, folgt aus
der gemessenen Lage des Bretts; das ist die \emph{Referenz}. Die Aufnahmen sind geteilt in \emph{Abstimmdaten} (zum
Einstellen) und \emph{Testdaten} (zum Prüfen); dazu kommt das Live-Log des Labortests vom 11.11.\ ohne Rohbilder.
Die \emph{Streuung} gibt an, wie stark die Spitze in einer Aufnahme um ihren Mittelwert zittert (RMS, das
quadratische Mittel der Abweichungen); ein gleichbleibender Abstand zur Referenz zählt nicht mit. Ein
\emph{Sprung} ist der Abstand der Spitze zwischen zwei aufeinanderfolgenden gezeichneten Bildern.

Von den vier Problemen erklärt (a) das feine Zittern, (b) und (c) die großen Sprünge und (d) die Abweichung
direkt nach dem Einrasten.

\subsection*{(a) Jedes Bild steht für sich}

Jedes Bild wird für sich ausgewertet; das Bild davor liefert nur den Startwert (Abschnitt~\ref{sec:heute}) und
einen Vergleich für Plausibilitätsprüfungen, beides aus seiner Farbpose. Das Rauschen jedes Bildes geht deshalb fast
ungefiltert auf die Spitze. Dazu kommt der
lange Hebel: Weil die Spitze 179,36\,mm vom Ursprung des Etiketts entfernt liegt, wird jeder kleine Kippfehler des
Etiketts an der Spitze vergrößert \cite[S.~535, 543]{west2004designing}. Gemessen: Mit eingespanntem Bohrer hat
die heute gezeichnete Spitze auf den Testdaten eine Streuung von 0,657\,mm. Steht der Bohrer still, könnte man über
viele Bilder mitteln und so das Zittern senken; das geschieht heute nicht. Ruhe ist am Patienten häufig, etwa mit
dem Bohrer am Eintrittspunkt.

\subsection*{(b) Umschalten zwischen zwei Schätzungen}

Welche der beiden Posen gezeichnet wird, entscheidet jedes Bild neu, und bei jedem Wechsel kann die Spitze um den
Unterschied der beiden Posen springen. Der carry-Versatz gleicht diesen Unterschied nur für das Bild aus, auf dem
er gemerkt wurde; liegt die Farbpose im nächsten Bild anders daneben, bleibt ein Rest. Im Labortest am 11.11.\
stand der Bohrer still (Bildbewegung der Ecken im Median 0,41\,px je Bild). Von den 2063 Bildern, für die die
Engine live eine neue Pose berechnet hatte, waren 286 (13,9\,\%) carry-Bilder, meist einzeln zwischen
fused-Bildern. Die Spitze sprang 83-mal um mehr als 5\,mm,
75-mal davon beim Wechsel zwischen fused und carry, höchstens um 12,84\,mm. Gemessen an der typischen Lage der
fused-Spitze (ihrem Median) lag die carry-Spitze im Median 3,67\,mm daneben, fast nur entlang der Sichtlinie der
Kamera, also genau in der Tiefe, in der die Farbpose schwach ist. Abbildung~\ref{fig:zeitreihe} zeigt dasselbe
Muster in einer Testaufnahme mit eingespanntem Bohrer: Die Ausschläge der heutigen Spitze fallen auf die
carry-Bilder.

\begin{figure}[!htb]
\centering
\begin{tikzpicture}
\begin{axis}[
    scale only axis, width=\linewidth-16mm, height=3.0cm,
    /pgf/number format/use comma,
    xlabel={Zeit [s]}, ylabel={Abstand zur Referenz [mm]},
    tick label style={font=\footnotesize}, label style={font=\footnotesize},
    xmin=0, xmax=6.1, ymin=0, ymax=7, ytick={0,1,2,3,4,5,6,7},
    unbounded coords=jump, enlarge x limits=0.01, clip mode=individual,
    axis line style={black!70}, tick style={black!70},
    ymajorgrids, grid style={black!12},
    legend style={font=\scriptsize, at={(0.5,1.03)}, anchor=south, legend columns=-1, draw=none,
                  /tikz/every even column/.append style={column sep=4mm}},
    /pgfplots/table/col sep=comma,
    /pgfplots/table/empty cells with={nan},
  ]
  \addplot[ColA, line width=0.6pt] table[x=t_s, y=err_A] {daten/zeitreihe_D05.csv};
  \addlegendentry{heute}
  \addplot[ColN, line width=1.2pt] table[x=t_s, y=err_neu] {daten/zeitreihe_D05.csv};
  \addlegendentry{neuer Schätzer}
  \addlegendimage{only marks, mark=*, mark size=1.4pt, draw=CarryC, fill=CarryC}
  \addlegendentry{heute, carry-Bild}
  \addplot[forget plot, only marks, scatter, scatter src=explicit symbolic,
    scatter/classes={fused={mark=none}, gate={mark=none}, rgb={mark=none},
                     carry={mark=*, mark size=1.4pt, draw=CarryC, fill=CarryC}}]
    table[x=t_s, y=err_A, meta=kind_A] {daten/zeitreihe_D05.csv};
\end{axis}
\end{tikzpicture}
\caption{Abstand der gezeichneten Spitze von der Referenz über eine ganze Aufnahme mit eingespanntem Bohrer, heute
(rot) und mit dem neuen Schätzer (grün). Orange Punkte: Bilder, auf denen heute die Farbpose plus dem gemerkten
Versatz gezeichnet wird (carry-Bilder). Dass beide Kurven nicht um null liegen, ist der bleibende Abstand der
Spitze zur Referenz. Beim ersten Bild fehlt die grüne Kurve, weil der neue Schätzer die Spitze direkt nach dem
ersten Finden des Etiketts noch nicht zeigt.}
\label{fig:zeitreihe}
\end{figure}

\subsection*{(c) Eine Prüfung verwirft die IR-Punkte gerade dann, wenn sie am meisten helfen würden}

Eine der vier Prüfungen verwirft die Fusion, wenn die fusionierte Pose den Ursprung des Etiketts mehr als 6\,mm
neben den der Farbpose legt. Sie misst damit den Fehler der Farbpose, nicht die Güte der IR-Messung: Liegt die
Farbpose weit daneben, muss die Fusion weit korrigieren und wird gerade deshalb verworfen. Im Labortest lagen auf
allen carry-Bildern 32 bis 50 IR-Punkte vor, die meist gut zum Modell passten (Restfehler im Median 0,81\,mm); auf
78,7\,\% der carry-Bilder blieb nur die 6-mm-Prüfung als Grund für die Ablehnung. Bei einer Ablehnung gehen alle
IR-Punkte des Bildes verloren. Üblich ist dagegen, auffällige Messungen stufenlos abzuwerten, statt sie mit harten
Schwellen zu verwerfen \cite[§9.3--9.4]{zhang1997parameter}.

\textbf{Wie (b) und (c) im Labortest zusammenwirkten.} Die Farbpose war dort schlecht bestimmt: Das Etikett lag
links oben am Bildrand, etwa 600\,px von der Bildmitte, unscharf und mit Farbsäumen. Die 6-mm-Prüfung sortierte
die Bilder. Auf Bildern, auf denen die Farbpose höchstens 6\,mm neben der fusionierten Pose lag, wurde fusioniert
und der Versatz gemerkt; dort lag der Ursprung des Etiketts in der Farbpose im Median 2,53\,mm neben dem der
fusionierten Pose. Auf Bildern, auf denen sie weiter danebenlag, wurde die Fusion verworfen und carry gezeichnet;
dort lag der Ursprung in der Farbpose im Median 6,8\,mm neben seiner typischen Lage in den fused-Bildern (ihrem
Median). Der gemerkte Versatz stammt also von einem Bild, auf dem die Farbpose besser war, und wird auf ein Bild
angewandt, auf dem sie schlechter ist. Deshalb sprang die Spitze beim Wechsel zwischen fused und carry.

\subsection*{(d) Unsicheres erstes Bild, keine Unsicherheit in der Anzeige}

Direkt nach dem Einrasten gibt es noch keinen gemerkten Versatz. Wird die Fusion auf diesem ersten Bild nicht
angenommen, zeichnet die Anzeige die reine Farbpose. Mit eingespanntem Bohrer trat die größte Abweichung der
heutigen Spitze vom Mittel ihrer Aufnahme in den Abstimm- wie in den Testdaten auf einem solchen ersten Bild mit
reiner Farbpose auf, in den Testdaten mit 6,07\,mm. Außerdem weiß die Anzeige nicht, wie sicher die Spitze ist;
eine Unsicherheit in Millimetern gibt es live nicht. Einziges Warnsignal ist Orange statt Grün im Kamerabild (in
der Röntgenansicht werden die Marken gedämpft), wenn weniger als 75\,\% des Etiketts sichtbar sind. Die Spitze
erscheint so, als wäre sie exakt; genau das kritisieren Simpson et al.\ an Navigationsanzeigen
\cite{simpson2014computation}.

%==========================================================================================
\section{Der Lösungsansatz}
\label{sec:neu}

Der neue Ansatz ersetzt das Umschalten durch \emph{einen} Schätzer, der als neue letzte Stufe hinter der heutigen
Posenbestimmung der Engine hinzukommt; Einstellungen und Posenbestimmung der Engine bleiben unverändert. Fachlich
ist er ein \emph{iterierter Kalman-Filter}. Ein Kalman-Filter verrechnet in jedem Bild eine Vorhersage aus den
früheren Bildern mit den neuen Messungen, jeweils danach, wie zuverlässig beide sind; \glqq iteriert\grqq{} heißt,
dass er seine Lösung in jedem Bild in drei kleinen Rechenschritten nachbessert. Abbildung~\ref{fig:filter} zeigt
den Kreislauf. Was ihn von heute unterscheidet; in Klammern steht, welches Problem aus Abschnitt~\ref{sec:problem}
der Punkt betrifft:

\begin{itemize}
  \item \textbf{Ein Gedächtnis (zu a).} Die Vorhersage für das neue Bild ist die Pose aus dem letzten Bild
        \cite[Bsp.~4.1, 4.2]{saerkkae2013bayesian}; ihre Unsicherheit wächst dabei je Bild, in Ruhe nur wenig. In
        Ruhe ist die Vorhersage deshalb nach einigen Bildern sicherer als ein einzelnes neues Bild, und jedes neue
        Bild verschiebt die Schätzung nur ein kleines Stück. Das wirkt wie ein Mittel über viele Bilder, in dem
        sich das Rauschen der einzelnen Bilder weitgehend ausgleicht. Dahinter steht die noch ungeprüfte Annahme,
        dass sehr schnelle freie Bewegungen am Patienten unwahrscheinlich sind.
  \item \textbf{Ein Ergebnis statt zwei (zu b).} Es gibt keinen Wechsel mehr zwischen zwei Schätzungen. Jede
        Messung zieht die Pose so weit, wie ihre Zuverlässigkeit es rechtfertigt; ein Bild mit IR-Punkten wird
        also nicht mehr nur ganz oder gar nicht verwendet: Gut passende IR-Punkte zählen voll, schlecht passende
        weniger. Ein Bild ohne IR-Punkte wird von der Vorhersage und den Farbecken getragen statt von einem
        eingefrorenen Versatz.
  \item \textbf{Alle Messungen, nach Zuverlässigkeit gewichtet.} In jedem Bild gehen alle Farbecken und, außer in
        den unten genannten Fällen, alle IR-Punkte ein, jede Messung mit ihrer eigenen Unsicherheit. Wie
        zuverlässig die beiden Messarten sind, wurde aus den beobachteten Abweichungen bei ruhender Spitze
        bestimmt.
  \item \textbf{Abwerten statt wegwerfen (zu c).} Je schlechter eine Messung zur gemeinsamen Lösung passt, desto
        geringer ist ihr Gewicht; so kann auch ein grober Ausreißer die Pose nur begrenzt ziehen
        \cite[§9.4]{zhang1997parameter}. Die 6-mm-Prüfung aus Abschnitt~\ref{sec:problem}\,(c), also der Abstand
        zwischen fusionierter Pose und Farbpose, entscheidet nicht mehr darüber, ob die IR-Punkte eines Bildes
        eingehen; gerade die Bilder, die heute an ihr scheitern, behalten so ihre IR-Punkte. Außen vor bleiben die
        IR-Punkte eines Bildes nur, wenn es weniger als 6 sind oder wenn sie grob nicht zum Modell passen. Diese
        Grenze ist keine neue Abstandsprüfung: Sie misst dieselbe Größe wie die heutige 2,5-mm-Prüfung aus
        Abschnitt~\ref{sec:heute}, nämlich wie gut die IR-Punkte zum Modell passen, lässt aber einen Restfehler bis
        6\,mm zu.
  \item \textbf{Folgen, wenn sich etwas bewegt.} Bewegen sich die Ecken des Etiketts im Bild gemeinsam, macht der
        Schätzer seine Vorhersage unsicherer und folgt dann vor allem den Messungen.
  \item \textbf{Schätzt seine Unsicherheit (zu d).} Mit der Pose führt er ihre Unsicherheit mit; daraus folgt die
        Unsicherheit der Spitze in Millimetern. Auch auf dem ersten Bild nach dem Einrasten verrechnet er
        Farbecken und IR-Punkte gemeinsam; eine Vorhersage gibt es dort noch nicht. Nach dem Einrasten bleibt die
        Spitze verborgen, bis ihre Unsicherheit höchstens 2\,mm beträgt, aber höchstens 5 Bilder lang (etwa
        167\,ms); in der geprüften Einstellung wird sie danach immer gezeichnet. Ob und wie die Unsicherheit
        selbst angezeigt wird, legt dieser Ansatz nicht fest.
\end{itemize}

\begin{figure}[!htb]
\centering
\begin{tikzpicture}[
  font=\footnotesize,
  st/.style={draw, rounded corners=1.5pt, align=left, text width=34mm, minimum height=27mm, inner sep=2.5pt,
             fill=white, execute at begin node={\hyphenpenalty=10000\exhyphenpenalty=10000}},
  inp/.style={draw, rounded corners=1.5pt, align=left, inner sep=2.5pt, fill=black!4,
              execute at begin node={\hyphenpenalty=10000\exhyphenpenalty=10000}},
  arr/.style={-{Stealth[length=1.8mm]}, semithick},
  fb/.style={-{Stealth[length=1.8mm]}, dashed, black!60},
  lab/.style={font=\footnotesize, fill=white, inner sep=1pt}
]
\node[st] (V) at (0,0) {\textbf{1 Vorhersage}\\[1pt] Pose wie zuvor; Unsicherheit wächst in Ruhe wenig, bei Bewegung
im Bild stark};
\node[st, right=5mm of V] (U) {\textbf{2 Abgleich}\\[1pt] Pose, die Vorhersage und Messungen am besten erklärt;
schlecht passende Messungen zählen weniger};
\node[st, right=5mm of U] (P) {\textbf{3 Neue Schätzung}\\[1pt] Pose mit Unsicherheit; daraus Spitze und
Unsicherheit in mm};
\node[st, right=5mm of P] (A) {\textbf{4 Anzeige}\\[1pt] Spitze zeichnen; nach dem Einrasten erst, wenn sie sicher
genug ist (höchstens 5 Bilder Wartezeit)};
\node[inp, text width=118mm, anchor=south west] (E) at ($(V.north west)+(0,9mm)$)
  {\textbf{Engine, neues Kamerabild:} Pose der Engine als Ausgangspunkt der Rechnung; Farbecken und IR-Punkte,
   jede Messung mit ihrer Unsicherheit};
\draw[arr] (V) -- (U);
\draw[arr] (U) -- (P);
\draw[arr] (P) -- (A);
\draw[arr] (E.south -| V) -- node[lab, right=0.5mm] {Bewegung im Bild} (V.north);
\draw[arr] (E.south -| U) -- node[lab, right=0.5mm] {Messungen, Ausgangspunkt} (U.north);
\draw[fb] (P.south) -- ++(0,-5mm) -| (V.south);
\node[lab, anchor=north] at ($(V.south)!0.5!(P.south)+(0,-5mm)$) {nächstes Bild: die neue Schätzung wird zur
Vorhersage};
\end{tikzpicture}
\caption{Ein Schritt des neuen Schätzers. Die Schätzung des letzten Bildes wird zur Vorhersage für das neue Bild;
der Abgleich verrechnet sie mit den Messungen des neuen Bildes, jeweils nach ihrer Zuverlässigkeit. Nach dem
Einrasten, also dem ersten Finden des Etiketts nach dem Start oder nach einem Verlust, gibt es noch keine
Vorhersage; dann zählen nur die Messungen.}
\label{fig:filter}
\end{figure}

Als Ausgangspunkt seiner Rechnung nimmt der Schätzer die Pose, die die Engine ausgibt: die fusionierte Pose, wenn
sie die Prüfungen besteht, sonst die Farbpose. Die Prüfungen entscheiden also nur noch, wo seine Rechnung beginnt;
den Startwert der Engine für die Farbpose (Abschnitt~\ref{sec:heute}) ändert er nicht. Bei nur drei
Rechenschritten beeinflusst der Ausgangspunkt das Ergebnis aber messbar, die Prüfungen wirken so schwach weiter.
Nach einer Lücke von
mindestens 0,5\,s oder wenn die Engine ihr Tracking zurücksetzt, beginnt der Schätzer ohne Gedächtnis neu.

\begin{mybox}{Für Interessierte: was der Schätzer in jedem Bild rechnet}
Zustand ist allein die Pose $\mathbf{x}_k = \Tcm$ (Etikett $\to$ Farbkamera) mit ihrer Kovarianz
$\boldsymbol{\Sigma}_k$. Die Vorhersage
lässt die Pose gleich und vergrößert nur ihre Unsicherheit um das Prozessrauschen $\mathbf{Q}_k$, klein in Ruhe
und groß bei Bewegung im Bild: $\xb = \xh_{k-1}$, $\Sb = \boldsymbol{\Sigma}_{k-1} + \mathbf{Q}_k$. Der Abgleich
sucht die Pose mit den kleinsten Kosten
\begin{equation*}
  \xh_k = \arg\min_{\mathbf{x}}\;
  \tfrac12\,\bigl\lVert \mathbf{x} \boxminus \xb \bigr\rVert^2_{\Sb^{-1}}
  + \sum_i \rho(e_i) + \sum_j \rho(e_j) .
\end{equation*}
Der erste Term ist die Abweichung von der Vorhersage, gewichtet mit deren Unsicherheit; $e_i$ und $e_j$ sind die
Abstände zwischen der aus der Pose $\mathbf{x}$ berechneten und der gemessenen Farbecke $i$ (im Bild) bzw.\ dem
IR-Punkt $j$ (im Raum), normiert mit der Unsicherheit der Messung. Die Huber-Funktion $\rho$ ist quadratisch bis
$c = 1{,}5$ und darüber linear, sodass große Abweichungen weniger zählen. Gelöst wird mit drei
Gauß-Newton-Schritten ab der Engine-Pose, die Gewichte jedes Mal neu berechnet. Mit rein quadratischen Kosten
statt $\rho$ ist das iterierte Kalman-Update ein Gauß-Newton-Verfahren auf solchen Kosten
\cite{bell1993iterated}, mit $\rho$ ein iterativ neu gewichtetes Kleinste-Quadrate-Verfahren
\cite[§9.4]{zhang1997parameter}. $\boldsymbol{\Sigma}_k$ ist die Inverse der Gauß-Newton-Näherung der Krümmung am
Ergebnis; aus ihr folgt die Unsicherheit $\stip$ der Spitze $\Ptc = \hat{\mathbf{R}}_k\,\Ptd + \hat{\mathbf{t}}_k$;
$\hat{\mathbf{R}}_k$ und $\hat{\mathbf{t}}_k$ sind Drehung und Verschiebung von $\xh_k$, $\Ptd$ ist die Lage der
Spitze im Werkzeugrahmen d, hier gleich dem Etikettrahmen m.
\end{mybox}

%==========================================================================================
\section{Was es bringt}
\label{sec:ergebnis}

\begin{table}[!htb]
\centering\small
\caption{Kernzahlen heute und mit dem neuen Schätzer. Streuung: RMS der Abweichung der Spitze vom Mittel ihrer
Aufnahme; Sprung: Abstand der Spitze zwischen zwei aufeinanderfolgenden gezeichneten Bildern; Nachlauf: wie viel
Zeit die Schätzung in Bewegung hinter einer Lösung zurückliegt, die nur das aktuelle Bild verwendet; bleibender
Abstand: Abstand zwischen dem Mittel der Spitze und der Referenz, gemittelt über die Aufnahmen. Der Schätzer wurde
vor der Auswertung der Testdaten eingefroren.}
\label{tab:kern}
\renewcommand{\arraystretch}{1.12}
\begin{tabularx}{\textwidth}{@{}>{\raggedright\arraybackslash}X>{\raggedright\arraybackslash}p{24mm}>{\raggedright\arraybackslash}p{56mm}@{}}
  \toprule
  Größe & heute & neuer Schätzer \\
  \midrule
  Streuung, Bohrer eingespannt (Testdaten) & 0,657\,mm & 0,215\,mm \\
  größte Abweichung vom Mittel, eingespannt (Testdaten) & 6,07\,mm & 1,07\,mm; ohne Verbergen nach dem Einrasten 2,07\,mm \\
  Streuung bei Verdeckung des Etiketts (Abstimmdaten; Testdaten ohne Verdeckung) & 1,344\,mm & 0,390\,mm \\
  Sprünge über 1\,mm, eingespannte Beispielaufnahme (Abb.~\ref{fig:zeitreihe}, Testdaten) & 29 & 0 \\
  Sprünge über 5\,mm, handgeführt, Etikett am Bildrand (Abstimmdaten) & 16 & 0 \\
  Nachlauf in Bewegung, je Testaufnahme & kein Nachlauf durch Glättung & höchstens 12,6\,ms; nur in einer der vier Aufnahmen ohne IR 166,6\,ms \\
  zusätzliche Rechenzeit je Bild & --- & 7--19\,µs (C++-Prototyp, Warmwert aus einer Messschleife; eingebaut
  mehr, Abschnitt~\ref{sec:cpp}) \\
  bleibender Abstand zur Referenz, eingespannt (Testdaten) & 1,80\,mm & 1,96\,mm \\
  \bottomrule
\end{tabularx}
\end{table}

\textbf{Einordnung.} Tabelle~\ref{tab:kern} fasst die Kernzahlen zusammen; alle Einstellungen des Schätzers wurden
vor dem Testlauf auf den Abstimmdaten festgelegt. Mit eingespanntem Bohrer ist die Spitze etwa dreimal ruhiger. Der
größte Ausschlag sinkt von
6,07 auf 1,07\,mm, teils durch das Verbergen nach dem Einrasten (ohne es 2,07\,mm), teils durch die Schätzung
selbst. Bei Verdeckung trägt die Vorhersage, und die verbliebenen Ecken korrigieren sie. In einer dem Labortest
ähnlichen Aufnahme verschwinden die großen Sprünge, die Streuung bleibt aber (Abschnitt~\ref{sec:grenzen}).

\textbf{Nachlauf und Rechenzeit.} Der Preis jeder Glättung ist, dass die Schätzung in Bewegung nachläuft. Die
heutige Ausgabe glättet
die Pose nicht und ist darum am schnellsten; der neue Schätzer läuft je Testaufnahme höchstens 12,6\,ms nach, mit
einer Ausnahme unter den Aufnahmen ohne IR (Abschnitt~\ref{sec:grenzen}). Die Rechenzeit fällt nicht ins Gewicht: 7
bis 19\,µs je Bild (Median je
Aufnahme, Warmwert des Prototyps aus einer Messschleife); die Engine selbst brauchte im Labortest 10,4\,ms (Median) bei 33\,ms Zeit je Bild.

%==========================================================================================
\section{Was es nicht löst}
\label{sec:grenzen}

\begin{itemize}
  \item \textbf{Der Abstand zur Referenz bleibt.} Die Spitze liegt im Mittel etwa 1,5 bis 2\,mm neben
        der Referenz, mit dem neuen Schätzer auf den Testdaten sogar etwas mehr als heute. Im Labortest lag sie
        gegenüber der Bohrerspitze, die aus dem AP- und dem Lateral-Röntgenbild bestimmt wurde, 2,6\,mm
        voraus entlang der Bohrerachse. Einen gleichbleibenden Messfehler erkennt ein Filter ohne zusätzliche
        Information nicht \cite[§5.1.3]{barfoot2017state}; das ist ein eigenes Thema.
  \item \textbf{Ohne IR ist die Tiefe schwach.} Von den vier Testaufnahmen ohne IR lief der Schätzer gegenüber der
        Lösung aus dem jeweiligen Bild allein in dreien höchstens 12,6\,ms nach, in einer 166,6\,ms. Dort bewegte
        sich das Etikett im Bild nur 0,1 bis 0,4\,px je Bild, während die Tiefe der Farbpose um 1 bis 7\,mm je
        Bild sprang; die Vorhersage öffnet sich aber nur bei Bewegung der Ecken im Bild. Ob sich die
        Spitze so bewegte, ist offen: Es gab keine Brett-Referenz, und die Tiefe der Einzelbildlösung
        springt selbst. Für die Nutzerstudie muss IR an sein.
  \item \textbf{Handgeführt am Bildrand nur teilweise besser.} In einer handgeführten Aufnahme mit dem Etikett am
        oberen Bildrand (Abstimmdaten), ähnlich wie im Labortest, verschwinden die großen Sprünge (über 5\,mm: von
        16 auf 0), und der größte Abstand zur Referenz sinkt von 11,93 auf 5,26\,mm. Die Streuung bleibt aber
        gleich (2,968 bzw.\ 2,970\,mm), und der Median des Abstands zur Referenz steigt aus offener Ursache von
        2,42 auf 3,06\,mm, vor allem am Anfang.
  \item \textbf{Die Unsicherheit ist eine Orientierung, keine Garantie.} Größere angegebene Unsicherheit geht mit
        größeren Fehlern einher, seltene große Fehler unterschätzt sie aber. Die Schwelle von 2\,mm für das erste
        Anzeigen nach dem Einrasten ist gesetzt, nicht validiert.
  \item \textbf{Neue Aufnahmen nötig.} Die Testdaten stammen aus denselben Sitzungen wie die Abstimmdaten. Nötig
  sind neue Rohaufnahmen mit dem Etikett am
        Bildrand wie im Labortest, mit langer Ruhe, mit bewegter Referenz und unter anderen Lichtverhältnissen.
\end{itemize}

%==========================================================================================
\section{Warum nicht einfacher, warum nicht aufwendiger?}
\label{sec:smoother}

\textbf{Einfacher: die letzte fused-Spitze halten.} Die einfachste Abhilfe wäre, auf carry-Bildern bei ruhigem
Bild die letzte fused-Spitze stehen zu lassen. Weil das Halten nur Größen braucht, die das Live-Log des Labortests
enthält, lässt es sich dort nachrechnen: Die Sprünge über 5\,mm sinken von 83 auf 0. Den neuen Schätzer kann man
dort nicht nachrechnen, weil er die Rohbilder braucht. Das Halten hilft aber nicht beim ersten Bild nach dem
Einrasten, nicht bei Verdeckung und nicht in Bewegung, und in einer Aufnahme mit Sprüngen auch zwischen zwei
fused-Bildern nur teilweise.

\textbf{Vorhanden: der Ausgabe-Filter der Engine.} Die Engine hat schon einen Kalman-Filter, der die fertige Pose
glättet. Er ist live aus: Er zog die Pose in Bewegung vom Bild weg und brachte auch nach einer Reparatur keinen
belastbaren Gewinn gegenüber carry. Er sieht nur die fertige Pose und kann
nicht erkennen, welche Ecken oder IR-Punkte gut sind. Der neue Schätzer arbeitet dagegen mit den einzelnen
Messungen; man spricht von enger Kopplung \cite[S.~493]{tobergte2009robust}.

\textbf{Aufwendiger: ein Glätter über ein Zeitfenster (Fixed-Lag-Smoother).} Ein solcher Glätter rechnet die Posen
eines kurzen Zeitfensters, etwa der letzten 0,5\,s, gemeinsam neu. Ältere Posen kann er verbessern, weil er für sie
auch die späteren Bilder kennt. Live gezeigt wird aber nur die \emph{neueste} Pose, und für sie gibt es noch kein
späteres Bild. Im idealisierten Fall (lineare Modelle, gaußsches Rauschen, in Glätter und Filter dieselben Modelle)
liefert er für sie genau dasselbe wie der Filter, sofern er die Information der Bilder vor dem Fenster nicht
verwirft, sondern weiterführt \cite[Thm.~8.2]{saerkkae2013bayesian}. Gemessen wurde das gegen eine frühere Fassung
des Filters, mit der der Glätter dasselbe Modell teilte: Auf den Abstimmdaten lagen die neuesten Spitzen beider im
Median 0,0044\,mm und in 99\,\% der Bilder höchstens 0,225\,mm auseinander. Der größte Unterschied, 24,52\,mm, trat
in schneller Bewegung in einer Aufnahme mit grob falschen IR-Punkten auf; solche IR-Punkte lässt der neue Schätzer
außen vor. Gewinnen könnte ein Glätter vor allem mit verzögerter Anzeige;
Fixed-Lag-Glättung ist Filtern mit einer tolerierten festen Verzögerung \cite[S.~160]{saerkkae2013bayesian}. Nach
einer einfachen Modellrechnung brächten 20 Bilder (0,67\,s) Verzögerung in Ruhe nur 20 bis 27\,\% weniger
Streuung. Für die Live-Anzeige am Patienten lohnt sich das nicht.

%==========================================================================================
\section{Stand der Umsetzung}
\label{sec:cpp}

Stand 13.11.2026: Der Schätzer ist in die C++-Engine eingebaut, an die Anzeige angeschlossen und in der
Live-Anzeige eingeschaltet. Wo die Abschnitte davor \glqq heute\grqq{} sagen, meinen sie den Stand vor dem Schätzer. In
der Engine sitzt er zwischen der IR-Fusion (Abb.~\ref{fig:heute}) und dem Ausgabe-Filter aus
Abschnitt~\ref{sec:smoother}. Er nimmt die Pose der Engine als Ausgangspunkt, bekommt von ihr die einzelnen Messungen
und schreibt seine Pose, die Spitze, deren Unsicherheit und die Angabe, ob gezeichnet wird, in eigene Ergebnisfelder.
Die Pose der Engine überschreibt er nicht: Sonst schlüge die bisherige Anzeige auf carry-Bildern ihren gemerkten
Versatz zusätzlich auf die Schätzung auf, und die Pose in Protokollen und Auswertungen bekäme eine andere Bedeutung. So
kann die bisherige Spitze zum Vergleich mitlaufen; im Evaluation-Modus der Anzeige wird sie mitprotokolliert.

\textbf{Eine Engine.} Es gibt nur noch eine Fassung der Engine: Die geprüfte neue Fassung hat am 13.11. die bisherige
ersetzt und wird ohne weitere Einstellung geladen. Die bisherige Fassung ist gesichert.

\textbf{Schalter.} Ein- und ausgeschaltet wird der Schätzer mit \texttt{LIVE\_DRILL\_ESTIMATOR}. Das ist eine
Einstellung in der Konfigurationsdatei des Trackings; sie steht auf True (Schätzer an) oder False (Schätzer aus).
Voreingestellt ist True. Die Einstellung gilt für die Live-Anzeige und für das Werkzeug, mit dem Testaufnahmen gemacht
werden; eine Änderung wirkt nach einem Neustart des Programms. Mit False sind Rechnung und Anzeige wie vor dem
Schätzer; das Evaluation-Log der Live-Anzeige hat aber das neue Format (Version 3), mit Vorgabewerten in den Spalten des Schätzers. Belegt
ist das auf acht Aufnahmen mit zusammen 2400 Bildern: Alle Ergebnisse bis auf die gemessenen Rechenzeiten sind
bitgleich zur bisherigen Engine, also bis aufs letzte Bit gleich, auch über die echte Live-Anzeige. Mit True bleiben
die bisherigen Ergebnisse der Engine ebenfalls bitgleich; der Schätzer kommt nur hinzu und wirkt nicht auf ihre
Posenbestimmung zurück. Dass der Vergleich selbst funktioniert, ist eigens geprüft: Er meldet schon eine absichtlich
eingebaute Änderung im letzten Bit. Die frühere Einstellung, die vor dem Programmstart gesetzt werden musste, gibt es
nicht mehr.

\textbf{Auswertungen der Masterarbeit.} Ihre Skripte legen die Einstellungen des Trackings selbst fest und rechnen
auch mit \texttt{LIVE\_DRILL\_ESTIMATOR} auf True ohne Schätzer. Auf den acht Aufnahmen rechnen mit der neuen Engine
alle 17 dort nachrechenbaren Einstellungen bitgleich zur bisherigen, aus Masterarbeit, Paper und der Live-Anzeige mit
ausgeschaltetem Schätzer.
Skripte mit einer eigenen Kopie der bisherigen Engine sind ohnehin nicht betroffen.

\textbf{Zurück zum alten Stand.} Für das Verhalten von vorher genügt \texttt{LIVE\_DRILL\_ESTIMATOR} auf False. Wer
die bisherige Engine-Fassung aus der Sicherung zurückholt, muss die Einstellung zusätzlich auf False setzen; sonst
meldet jeder Start, dass der Schätzer aus bleibt (die Live-Anzeige rot), und gerechnet wird wie mit ausgeschaltetem
Schätzer.

\textbf{Ausgabe-Filter.} Ist der Schätzer eingeschaltet, schaltet die Engine den Ausgabe-Filter zwingend ab, damit nur
der Schätzer glättet. Ohne Schätzer bleibt der Filter wie bisher schaltbar, weil Auswertungen der Masterarbeit ihn
verwenden, und er rechnet dann genau wie bisher; live ist er ohnehin aus. Beides ist auf den acht Aufnahmen geprüft.

\textbf{Anzeige.} Bei eingeschaltetem Schätzer zeichnet die Anzeige seine Spitze, ohne carry; die vier Prüfungen
bestimmen nur noch, wo seine Rechnung beginnt. Ist die Spitze verborgen, verhält sich die Anzeige wie bei einem Bild
ohne Pose. Beim Start prüft das Programm, ob die geladene Engine den Schätzer kennt und richtig eingestellt ist; wenn
nicht, läuft das Tracking wie vorher, mit einer roten Meldung. In der Ansicht mit der schmalen Bedienleiste am Rand,
die keine Statuszeile hat, steht diese Meldung als roter Vermerk vor dem Status, eine Warnung als gelber; die ganze
Meldung erscheint, wenn man mit der Maus darauf zeigt. Der Fehler, dass dieser Status nach dem Stoppen rot oder gelb
blieb, ist behoben. Auf den acht Aufnahmen liefert der eingebaute Schätzer dieselben Spitzen wie die geprüfte
Python-Fassung, bis auf Rundungsunterschiede unter $10^{-12}$\,mm.

\textbf{Fehlerzähler.} Beim Stoppen schreibt das Programm bei eingeschaltetem Schätzer eine Zeile mit fünf
Fehlerzählern ins Protokoll, etwa wie oft der Schätzer keine brauchbare Pose lieferte. Erwartet ist überall null; in
allen Prüfungen ohne absichtlich ausgelöste Fehler war es so. Zwei der Zähler, Fehler in der Schätzerstufe der Engine
und beim Speichern der
Schätzer-Eingabe, zählen seit dem Aufbau des Trackers und nicht je Lauf, denn das Programm behält den Tracker über
mehrere Läufe.

\textbf{Labortest.} Er läuft im Evaluation-Modus der Live-Anzeige und wird so ausgewertet, wie er angezeigt wurde.
Dieser Modus speichert keine Rohbilder. Nachrechnen lässt sich der Labortest also nicht, und die neuen Rohaufnahmen aus
Abschnitt~\ref{sec:grenzen} entstehen dabei nicht.

\textbf{Rechenzeit.} Gemessen beim Nachrechnen der acht Aufnahmen: In jedem Messlauf brauchte die ganze Stufe in 99 von
100 Bildern höchstens 0,211\,ms; erlaubt sind 0,33\,ms, also 1\,\% der 33\,ms, die für ein Bild zur Verfügung stehen.
Im Median sind es etwa 0,07\,ms je Bild, mehr als beim Prototyp in Tabelle~\ref{tab:kern}, weil die Stufe die
Messungen erst aufbereitet. Das Umwandeln der neuen Ergebnisse auf der Python-Seite dauert etwa 20\,µs je Bild. Über
alle Bilder gewertet hält es seine Grenze von 20\,µs knapp; für Bilder mit Schätzung lag es in einer Messreihe
0,2\,µs darüber. Neben den 33\,ms je Bild fällt das nicht ins Gewicht. Im laufenden Betrieb ist die Rechenzeit noch
nicht gemessen.

\textbf{Zusätze live.} Live laufen drei Zusätze mit, die von der geprüften Einstellung abweichen: Der Schätzer
beginnt schon nach einer Lücke von 0,2\,s statt 0,5\,s neu, ebenso, wenn die Unsicherheit, die seine Vorhersage in
einem Bild hinzubekommt, 20\,mm übersteigt (bei starker oder unbekannter Bewegung im Bild), und die Spitze bleibt
verborgen, solange ihre Unsicherheit über 5\,mm liegt. Festgelegt wurden die Zusätze auf den Abstimmdaten und an
einem nachgestellten Ausfall; dort verringerten sie vor allem die größten Abweichungen. Auf den Testdaten, die danach
einmal ohne Änderung eines Werts ausgewertet wurden, bewirken sie mit der Spitze im Sitz nichts und in Bewegung kaum
etwas; dort bleibt die Spitze dafür in 8,8 statt 2,2\,\% der Bilder verborgen.

\textbf{Was noch aussteht.} Noch aus steht der Labortest; mit echter Kamera sind die Anzeige und die Rechenzeit bisher
nicht geprüft. Offen sind außerdem einige kleinere Entscheidungen, etwa zusätzliche Prüfungen der Eingaben in der
Engine und ein Test der Anzeige mit der echten Engine. In die Versionsverwaltung übernommen (committet) ist noch
nichts; das geschieht nur auf Auftrag.

\begin{otherlanguage}{english}
\begin{multicols}{2}[\section*{Literatur}]
\raggedcolumns
\footnotesize
\raggedright
\bibliographystyle{unsrt}
\bibliography{Tracking_Schaetzer}
\end{multicols}
\end{otherlanguage}

\end{document}
